\chapter{Use Cases}

This appendix describes the primary use cases of the multi-camera tracking system as exercised through the web application. Each use case maps to one step of the operator workflow.

\begin{figure}[htbp]
\centering
\framebox[0.85\textwidth]{\parbox{0.83\textwidth}{\centering\vspace{2.5cm}\textit{Figure placeholder: UML use case diagram (Actor: Operator, UC1--UC7)}\vspace{2.5cm}}}
% TODO: insert figure — UML use case diagram with actor "Operator" and use cases UC1–UC7
\caption{System use case diagram.}
\label{fig:use_case}
\end{figure}

\section*{UC1: Upload Camera Footage}
\textbf{Actor:} Operator \\
\textbf{Precondition:} Operator has access to the web application and has one or more video files from surveillance cameras. \\
\textbf{Main Flow:}
\begin{enumerate}
    \item Operator navigates to the Upload screen.
    \item Operator drags and drops one or more video files onto the upload zone, or uses the file picker.
    \item System uploads files to the backend, creates video catalogue entries, and displays a thumbnail gallery.
    \item System automatically triggers Stage~0 (frame extraction) and Stage~1 (detection) on the uploaded videos.
\end{enumerate}
\textbf{Postcondition:} Video files are stored; per-frame detections are available.

\section*{UC2: Review Detections}
\textbf{Actor:} Operator \\
\textbf{Precondition:} UC1 complete; Stage~1 detection results available. \\
\textbf{Main Flow:}
\begin{enumerate}
    \item System renders YOLO bounding boxes overlaid on the video player.
    \item Operator reviews detections frame by frame or via video scrubbing.
    \item Operator verifies that vehicles of interest are correctly detected.
\end{enumerate}
\textbf{Postcondition:} Operator has confirmed detection quality.

\section*{UC3: Select Target of Interest}
\textbf{Actor:} Operator \\
\textbf{Precondition:} Detections displayed (UC2 complete). \\
\textbf{Main Flow:}
\begin{enumerate}
    \item Operator clicks on a bounding box in the video frame.
    \item System highlights the selected box in green and records the associated tracklet ID.
    \item Operator may select multiple targets or deselect by clicking again.
\end{enumerate}
\textbf{Postcondition:} One or more tracklets are designated as query targets.

\section*{UC4: Run Cross-Camera Inference}
\textbf{Actor:} Operator \\
\textbf{Precondition:} At least one target selected (UC3 complete). \\
\textbf{Main Flow:}
\begin{enumerate}
    \item Operator optionally fills in location metadata (country, governorate, district, intersection) and a datetime range.
    \item Operator clicks ``Run Inference''.
    \item System executes Stages~2--4: feature extraction, FAISS indexing, cross-camera association.
    \item System assigns global track IDs and returns results.
\end{enumerate}
\textbf{Postcondition:} Global trajectories are available.

\section*{UC5: Review Timeline}
\textbf{Actor:} Operator \\
\textbf{Precondition:} Cross-camera results available (UC4 complete). \\
\textbf{Main Flow:}
\begin{enumerate}
    \item System displays a Clipchamp-style timeline with tracklet segments colour-coded by camera.
    \item A split-screen player shows the video for the selected camera segment.
    \item Operator reviews the sequence of camera appearances for each global identity.
    \item Operator may drag segment boundaries or reorder segments.
\end{enumerate}
\textbf{Postcondition:} Operator has reviewed the cross-camera trajectory sequence.

\section*{UC6: Refine Search}
\textbf{Actor:} Operator \\
\textbf{Precondition:} Initial results reviewed (UC5 complete). \\
\textbf{Main Flow:}
\begin{enumerate}
    \item Operator selects specific reference frames from the timeline.
    \item Operator adjusts video playback speed (e.g., 0.25$\times$ for close inspection).
    \item Operator triggers a re-search using the selected frames as new query embeddings.
    \item System updates results with the refined query.
\end{enumerate}
\textbf{Postcondition:} Refined global trajectory returned.

\section*{UC7: Export Results}
\textbf{Actor:} Operator \\
\textbf{Precondition:} Final results reviewed. \\
\textbf{Main Flow:}
\begin{enumerate}
    \item Operator selects the desired export format: summarised video, MOT CSV, or JSON trajectory.
    \item Operator optionally selects a camera grid layout (1$\times$1 to 5$\times$5).
    \item System generates and downloads the export file.
\end{enumerate}
\textbf{Postcondition:} Export file saved to operator's local machine.
